% TEMPLATE for Usenix papers, specifically to meet requirements of
%  USENIX '05
% originally a template for producing IEEE-format articles using LaTeX.
%   written by Matthew Ward, CS Department, Worcester Polytechnic Institute.
% adapted by David Beazley for his excellent SWIG paper in Proceedings,
%   Tcl 96
% turned into a smartass generic template by De Clarke, with thanks to
%   both the above pioneers
% use at your own risk.  Complaints to /dev/null.
% make it two column with no page numbering, default is 10 point

% Munged by Fred Douglis <douglis@research.att.com> 10/97 to separate
% the .sty file from the LaTeX source template, so that people can
% more easily include the .sty file into an existing document.  Also
% changed to more closely follow the style guidelines as represented
% by the Word sample file. 

% Note that since 2010, USENIX does not require endnotes. If you want
% foot of page notes, don't include the endnotes package in the 
% usepackage command, below.

\documentclass[letterpaper,twocolumn,10pt]{article}
\usepackage{usenix,epsfig,endnotes}
\begin{document}

%don't want date printed
\date{}

%make title bold and 14 pt font (Latex default is non-bold, 16 pt)
\title{\Large \bf C.H.U.D. : Catastrophic Harmful Update Detection in Sequential Fine-Tuning}

\author{
{\rm Jorge Manuel Torre}\\
Virginia Tech
\and
{\rm Elias Mengistu}\\
Virginia Tech  \\
\and
{\rm Alex Levy}\\
Virginia Tech \\
\and
{\rm Jeevan Shahi}\\
Virginia Tech
}

\maketitle

% Use the following at camera-ready time to suppress page numbers.
% Comment it out when you first submit the paper for review.
\thispagestyle{empty}


\subsection*{Abstract}

% The safety alignment in modern open-weight large language models is usually achieved through techniques such as supervised fine-tuning and reinforcement learning from human feedback. However there's been a recent trend in which research suggests that these safety behaviors only occupy a part of the parameter space and can degrade under additional repeated fine-tuning. In this work, we investigate whether catastrophic forgetting can be leveraged as an attack mechanism against models that are safety-aligned. We propose a two-stage sequential fine-tuning attack in which an attacker first performs a benign fine-tuning phase to erode safety representations and then a harmful data injection phase to change model behavior. 

% To study this threat, we have decided to evaluate a sequential attack pipeline using GSM8K as the benign dataset that we use and BeaverTails as the harmful dataset in the second phase. We measure the safety degradation using attack success rate and then test whether the LoX (Low-Rank Extrapolation) defense remains sturdy under this attack strategy. Our evaluation compares the baseline, the attacked, and the LoX-protected models to fully determine whether LoX can actually preserve safety alignment when catastrophic forgetting is used as a initial step. 

% Our goal is to change the mindset of catastrophic forgetting from an issue created into something that has actual practical cases and provide a better and more rigorous benchmark for evaluating these defenses against fine-tuning attacks on already deployed systems.

Safety alignment in open-weight models is increasingly understood to be a fragile, low-rank modification layered on top of the models' base parameters. This alignment has been shown to degrade under continued fine-tuning, even on entirely benign data. We deliberately weaponize this vulnerability as an attack mechanism to stress-test LoX (Low-Rank Extrapolation), a recently proposed defense against fine-tuning attacks \cite{perin2025loxlowrankextrapolationrobustifies}. Our approach, C.H.U.D. (Catastrophic Harmful Update Detection), is a two-stage sequential fine-tuning attack in which an adversary first exploits catastrophic forgetting to quietly erode a model's safety alignment through benign fine-tuning, before introducing harmful data to redirect the model's behavior. This introduces a threat model that LoX was not explicitly designed to handle.

Using Llama-2-7B-Chat as our target, we initiate this attack by fine-tuning on GSM8K, a math reasoning dataset with no safety-relevant content, to induce catastrophic forgetting of alignment properties \cite{cobbe2021training}. Following this, we fine-tune on BeaverTails, a dataset containing harmful samples introduced in \cite{ji2023beavertails}. We measure safety degradation via Attack Success Rate (ASR) across three conditions: the unaligned baseline, the sequential attack-model, and a model hardened with LoX. By reframing catastrophic forgetting from an incidental training artifact into an active attack, this work seeks to establish a more rigorous benchmark for evaluating the robustness of fine-tuning defenses in deployed open-weight models. 

\section{Motivation \& Problem Statement}

\subsection{Motivation}

The democratization of Large Language Models (LLMs) through 
open-weight releases has introduced a critical security paradox. 
While models like Llama-2-7B-Instruct are released with extensive 
safety alignment, this protection is not an immutable property but 
rather a shallow, low-rank modification layered atop the model's 
base parameters~\cite{perin2025loxlowrankextrapolationrobustifies}. Recent research has demonstrated that 
this alignment is highly susceptible to drift during subsequent 
fine-tuning --- even when the training corpus is entirely 
benign~\cite{li-etal-2024-revisiting,huang2024harmfulfinetuningattacksdefenses}. This phenomenon, known as catastrophic 
forgetting, systematically degrades a model's safety subspace 
without any explicit adversarial signal, making it both difficult 
to detect and straightforward to exploit.

\subsection{Problem Statement}

While safety alignment provides an essential guardrail for deployed LLMs, its tendency to fail under continued fine-tuning demands equally robust defensive countermeasures \cite{cui2025persistentbackdoorattackscontinual, huang2024harmfulfinetuningattacksdefenses}. One recently proposed solution is LoX (Low-Rank Extrapolation), a post-alignment stage defensive technique that identifies and amplifies safety-critical low-rank directions in a models' parameter space, increasing its resistance to fine-tuning attacks before deployment \cite{perin2025loxlowrankextrapolationrobustifies}. Evaluations of LoX have demonstrated meaningful improvements in robustness against direct harmful fine-tuning. However, these evaluations consider only single-stage adversarial attacks in which the attacker immediately attempts fine-tuning to optimize harmful outputs.

This work investigates whether LoX's robustness guarantees hold under a more subtle and realistic threat: a two-stage sequential fine-tuning attack that exploits catastrophic forgetting as a mechanism for weakening the target. Rather than targeting safety alignment directly, an adversary first fine-tunes the model on entirely innocuous data, silently eroding its alignment properties before introducing any harmful inputs. By the time harmful fine-tuning begins in the second stage, the models' safety subspace has already been quietly undermined. 

To model this threat, we use GSM8K as our benign fine-tuning corpus and the BeaverTails dataset as our harmful fine-tuning source, evaluating the full attack pipeline against both an unprotected baseline (Llama-2-7b-chat) and a LoX-hardened model \cite{cobbe2021training, ji2023beavertails}. Our goal is to determine whether catastrophic forgetting represents a blind spot in LoX's threat model. In doing so, we hope to establish a more robust benchmark for alignment defenses in a broader scope.

We hypothesize that LoX's robustness will degrade under sequential fine-tuning because its amplification of safety-aligned directions does not prevent their gradual erosion under sufficient benign updates. As a result, models hardened with LoX will exhibit significantly higher ASR scores after the two-stage attack than under single-stage evaluations.

\subsection{Threat Model}

\paragraph{Attacker Model:} We model an adversary with complete access to the weights of a publicly released, safety-aligned open-weight model. The attacker possesses the computation resources necessary to perform parameter-efficient fine-tuning (e.g., via LoRA-based adaptation) and has access to two datasets: a benign domain-specific corpus used in the first attack stage, and a curated harmful dataset used in the second. For the benign stage, we use GSM8K: a dataset of 8,500 grade-school mathematics word problems requiring multi-step arithmetic reasoning \cite{cobbe2021training}. GSM8K was selected deliberately for its complete semantic distance from any safety-relevant content domain; mathematical reasoning problems contain no harmful signal, no ambiguous content, and bear no relationship to the safety behaviors the model was aligned to preserve. This makes it an ideal instrument for isolating catastrophic forgetting as the sole mechanism of safety degradation.

The attacker's primary objective is to produce a model that complies with harmful prompts that the original safety-aligned model would refuse, measured concretely by attack success rate (ASR) on a harmful evaluation set. Critically, the attacker is not assumed to have any knowledge of the internal structure of the LoX defense or the safety-critical subspaces it protects.


\paragraph{Defender Model:} The defender here is the model provider. Prior to release, the defender applies the LoX hardening procedure to the safety-aligned model, modifying its weights to increase robustness against fine-tuning attacks. After release, the defender has no visibility into how downstream parties use the model (i.e., they cannot monitor or constrain future fine-tuning pipelines). This assumption reflects the realistic constraints of the open-weight deployment mode, where the provider relinquishes complete control of the model upon release.

The defender's goal is to ensure that the model's safety alignment is preserved even after arbitrary downstream fine-tuning. In particular, a successful defense should prevent the product of policy-violating outputs, regardless of whether the attacker targets alignment directly or erodes it indirectly through catastrophic forgetting.

\section{Related Work}

\subsection{Safety Defenses in Open-Weight Models}

Modern open-weight LLMs are typically safety-aligned using supervised fine-tuning (SFT), reinforcement learning from human feedback (RLHF), or related preference-optimization techniques such as Direct Preference Optimization (DPO). These approaches modify a pretrained base model to reduce harmful or policy-violating outputs while preserving general capability. However, a growing body of work has shown that such alignment is not deeply integrated into model representations; rather, it often occupies a relatively small and fragile region of parameter space that can be disrupted by modest gradient updates \cite{li2025salorasafetyalignmentpreservedlowrank}.
After the risks of fine-tuning-based alignment degradation were identified, several categories of mitigation strategies emerged. One line of work modifies the alignment stage itself to account for robustness against future fine-tuning attacks \cite{rosati2024representationnoisingdefencemechanism}. Another targets the fine-tuning stage directly, constraining updates to mitigate safety degradation; for example, by retaining base model features or incorporating adversarial training \cite{wang2024mitigatingfinetuningbasedjailbreak}. A third category focuses on post-release recovery, restoring safety properties in models that have already been fine-tuned \cite{dong2021pretrainedlanguagemodelsfinetuned}.

LoX (Low-Rank Extrapolation) occupies a distinct position among these defenses. Rather than modifying the alignment or fine-tuning pipeline, LoX operates purely at the post-alignment stage, requiring only access to aligned and unaligned model checkpoints. It identifies safety-critical low-rank directions in the difference between these checkpoints and extrapolates along them to amplify safety-relevant components, increasing robustness to downstream fine-tuning without any changes to training or inference. This makes it broadly applicable regardless of how the model was aligned or how it will subsequently be fine-tuned. Building on earlier work showing that alignment can be enhanced by extrapolating between initial and aligned model weights (ExPO), LoX extends this idea by selectively targeting the safety subspace rather than extrapolating globally \cite{zheng2025modelextrapolationexpeditesalignment}.

Our work directly stress-tests LoX within this scope. Because LoX intervenes only at the post-alignment stage and makes no assumptions about the downstream fine-tuning process, it is in principle exposed to any attack that can erode safety representations after deployment, including the catastrophic forgetting-based preparatory stage we propose. Whether its low-rank extrapolation is sufficient to withstand indirect alignment erosion, rather than direct adversarial fine-tuning, is precisely the question this work investigates.

\subsection{Fine-Tuning Attacks on Aligned Models}

A related line of work examines the vulnerability of aligned LLMs to post-deployment fine-tuning attacks. These attacks typically assume that an adversary obtains access to model weights and performs additional fine-tuning to induce harmful behavior. 

The most straightforward attack strategy involves direct harmful fine-tuning, in which an adversary trains the model on an explicitly adversarial dataset designed to elicit policy-violating outputs. This approach was demonstrated to be remarkably effective even at small scales; just a few hundred harmful samples are sufficient to substantially degrade alignment in models that underwent extensive safety training \cite{huang2024harmfulfinetuningattacksdefenses}. Importantly, this work showed that even fine-tuning on benign, domain-specific tasks can unintentionally undermine safety-alignment. 

A second class of attacks focuses on backdoor insertion during supervised fine-tuning. Here, the adversary embeds hidden triggers into the training data such that the model behaves safely under normal conditions but produces harmful outputs when a specific trigger is present. These attacks are particularly concerning because the backdoor can persist through subsequent fine-tuning rounds, making them difficult to detect or remove after the fact \cite{cui2025persistentbackdoorattackscontinual}.

A third approach exploits low-rank adaptation methods such as LoRA to minimize computational cost while achieving substantial behavioral shifts \cite{ran2025loraleakmembershipinferenceattacks}. Because LoRA confines updates to a low-dimensional subspace of the parameter space, an attacker can redirect model behavior with far fewer resources than full fine-tuning requires, which lowers the barrier to entry for adversaries.

However, most existing evaluations focus on single-stage attacks in which the adversary directly optimizes towards harmful outputs. Our work extends this threat model by explicitly evaluating a two-phase attack: benign fine-tuning to erode safety representations through catastrophic forgetting, followed by harmful fine-tuning to redirect behavior. This exposes a gap in how existing defenses, including LoX, have been evaluated.

\subsection{Catastrophic Forgetting in Large Language Models}

Catastrophic forgetting has long been studied in continual learning and neural network research. In the context of LLMs, it has been observed that fine-tuning on domain-specific or task-specific data can degrade performance on unrelated tasks\cite{doi:10.1073/pnas.1611835114}. More recently, researchers have identified safety alignment as particularly vulnerable to such degradation\cite{alssum2025unforgottensafetypreservingsafety}.

Unlike traditional continual learning settings, safety alignment may correspond to relatively low-rank or low-magnitude parameter modifications layered atop a powerful pretrained base model. Consequently, modest gradient updates can disproportionately affect safety-related subspaces. While prior work has documented this drift, systematic evaluation of defenses under catastrophic-forgetting-based attack preparation remains limited \cite{ran2025loraleakmembershipinferenceattacks}.

Our contribution lies in explicitly weaponizing catastrophic forgetting as a preparatory mechanism and testing whether existing defenses remain robust under this more subtle threat model.

\subsection{Low-Rank Defenses and the LoX Framework}

Low-rank parameterization techniques, such as LoRA, have been widely used to reduce the computational cost of fine-tuning, as in the recent SaLoRA defense\cite{li2025salorasafetyalignmentpreservedlowrank}. Prior work, like SaLoRA leverages the observation that alignment updates often exhibit a low-rank structure, and propose constraining or regularizing these updates to preserve safety-relevant directions in parameter space.

Building off of this perspective, the Low-Rank Extrapolation (LoX) framework introduces a training-free defense that directly amplifies alignment-relevant components of a model. Rather than modifying the fine-tuning procedure, LoX operators post-hoc by decomposing the alingment update, $\Delta W_{\mbox{align}}$, into its principal components (using the singular value decomposition) and extrapolating along its top-ranked directions. These directions correspond to the most salient safety-related subspaces learned during alignment. By reinforcing them, LoX increases the model's resistance to subsequent fine-tuning that might otherwise degrade safety behavior. 

Formally, LoX constructs an updated parameter set by augmenting the aligned model with a scaled project of $\Delta W_{\mbox{align}}$ onto its top-$k$ singular components, where $k$ is chosen as the minimum rank required to preserve safety performance under an Attack Success Rate criterion. This formulation isolates and strengthens the most safety-critical structure in the alignment update while avoiding unnecessary parameter modifications.

Empirical results reported in the LoX work show that this approach improves robustness against direct adversarial fine-tuning compared to unprotected models. However, these evaluations are limited to single-stage attack settings.

\subsection{Positioning of This Work}

Our work occupies a gap at the intersection of catastrophic forgetting and fine-tuning defenses. Specifically:

\begin{itemize}
    \item Prior work shows that alignment is fragile under fine-tuning.
    \item Fine-tuning attack literature demonstrates that harmful behavior can be induced post hoc.
    \item The LoX defense strengthens robustness against direct fine-tuning attacks.
\end{itemize}

We integrate these threads by asking a new question: \emph{Does LoX remain effective when the attacker first degrades safety through benign fine-tuning before introducing harmful data?}

By stress-testing LoX under a sequential catastrophic-forgetting threat model, we evaluate whether low-rank extrapolation defenses protect against indirect erosion of alignment, not merely explicit adversarial optimization. This reframes catastrophic forgetting from an incidental training pathology into an active adversarial tool and provides a more stringent robustness benchmark for future alignment defenses.


\section{Proposed Work}

Our methodology proceeds in three phases:
\begin{enumerate}
    \item[1.] Establishing baseline safety and utility metrics on an unmodified model.
    \item[2.] Executing a two-stage sequential fine-tuning attack on a safety-aligned model without any defense, and:
    \item[3.] Applying the LoX defense and repeating the identical attack under the same conditions.
\end{enumerate}
All experiments use Llama-2-7B-Chat as the target model, consistent with the original LoX evaluation, allowing us to isolate the effect of our sequential attack strategy.

\subsection{Phase 1: Baseline Evaluation}

Before any fine-tuning, we evaluate the model's safety alignment and general utility to establish reference metrics for the later results. Safety is measured with Attack Success Rate (ASR), using Llama-3-8B-Guard as the judge model with an adapted evaluation prompt from the original LoX paper, based on Meta's Llama Usage Policy. Each response is flagged as safe or unsafe, and ASR is computed as the proportion of responses judged unsafe. General utility is measured using the utility score evaluation protocol from the LoX paper, ensuring that subsequent fine-tuning stages can be assessed not only for safety degradation but also for any unintended loss of capability.

\subsection{Phase 2: Sequential Fine-Tuning Attack}

The attack proceeds in two sequential stages, both using Low-Rank Adaptation (LoRA) for parameter-efficient fine-tuning:

\begin{itemize}
    \item \textbf{Stage 1 --- Benign Fine-Tuning:} We fine-tune the baseline model on GSM8K, a corpus of grade-school mathematics word problems with no safety-relevant content \cite{cobbe2021training}. The complete semantic distance between GSM8K and any safety-aligned behavior ensures that any observed degradation in alignment after this stage can be attributed solely to catastrophic forgetting rather than to any harmful signal in the training data. 
    
    \item \textbf{Stage 2 --- Harmful Fine-Tuning:} We continue fine-tuning the Stage 1 model on a subset of the BeaverTails dataset, which contains prompts and responses spanning multiple unsafe categories including hate speech, violence, and illegal activity~\cite{ji2023beavertails}. We evaluate ASR and utility scores after this stage, and use the difference between post-stage 2 metrics and baseline ones to quantify the full effect of the sequential attack.
\end{itemize}

To study how attack effectiveness scales with data, we'll run both fine-tuning stages across increasing sample sizes (10, 50, 100) from the respective datasets. This allows us to characterize how much data an adversary needs at each stage to meaninfully degrade safety alignment. 

\subsection{Phase 3: Defense Evaluation}

We apply the LoX defense~\cite{perin2025loxlowrankextrapolationrobustifies} to the original unmodified baseline model, following the procedure described in the original paper. We then repeat the identical two-stage attack from phase 2 on the LoX-protected model.

Post-attack metrics from the protected model are compared against both the unprotected baseline and the unprotected attacked model. This three way comparison allows us to isolate how much damage the sequential attack causes without any defense, and how much of that harm LoX successfully mitigates.

\subsection{Datasets and Models}

\paragraph{Target Models} We use Llama-2-7B-Chat~\cite{touvron2023llama} as our target model. Llama-2-7B-Chat was selected due to the original LoX evaluation including Llama-2-7B as a primary architecture~\cite{perin2025loxlowrankextrapolationrobustifies}, allowing us to isolate the effect of our sequential attack strategy rather than introducing confounds from an untested model.

\paragraph{Benign Dataset (Stage 1).} We use GSM8K~\cite{cobbe2021training}, a corpus of approximately 8,500 grade-school mathematics word problems requiring multi-step arithmetic reasoning. As discussed in Section 1.3, GSM8K was chosen for its complete semantic distance from any safety-relevant content, ensuring that any observed safety degradation after Stage 1 can be attributed solely to catastrophic forgetting rather than to harmful signal in the training data.

\paragraph{Harmful Dataset (Stage 2).} We use the BeaverTails dataset~\cite{ji2023beavertails}, which covers a broad range of unsafe content categories. BeaverTails is used exclusively in Stage 2; safety evaluation is conducted separately using AdvBench to ensure a clean separation between training/evaluation, and to more closely match the original evaluation in the LoX paper.

\paragraph{Safety Evaluation.} We measure Attack Success Rate (ASR) using the first 100 examples from AdvBench~\cite{zou2023universal}, using Llama-3-8B-Guard as the judge. A response is counted toward ASR if the judge deems it unsafe.

\paragraph{General Capability Benchmark.} We report the utility score as defined in the original LoX evaluation to confirm that neither stage of fine-tuning causes unintended degradation of general model capability alongside the intended safety degradation in~\cite{perin2025loxlowrankextrapolationrobustifies}.

\paragraph{Fine-Tuning Configuration.} Both stages of fine-tuning use Low-Rank Adaptation (LoRA)~\cite{hu2022lora}. We adopt the following hyperparameters, aligned with the configurations reported in the LoX paper where applicable: a LoRA rank of 8, learning rate of $5 \times 10^{-5}$, batch size of 20, and training for 2 epochs in Stage 1 (benign) and 5 epochs in Stage 2 (harmful). The maximum sequence length is set to 1024 tokens. These settings may be adjusted during experimentation; final values will be reported with our results.

\section{Limitations \& Future Work}
\label{sec:limitations}

\paragraph{Limitations.}
Our study is a focused stress-test of the LoX defense under a two-stage sequential fine-tuning attack. This limited scope raises a few limitations.

First, we evaluate only a single defense mechanism. While LoX represents a compelling and recently proposed post-alignment hardening technique, catastrophic forgettings' interaction with the other defenses mentioned in this proposal, such as alignment-stage or post-fine-tuning stage defenses, remains unexplored. Conclusions about LoX's robustness may or may not generalize to these alternative approaches.

Second, our results are specific to the Llama-2-7B-Chat model. Findings on this architecture may not transfer directly to others. In particular, larger models may exhibit different responses to catastrophic forgetting, and models aligned using different procedures may respond differently to benign fine-tuning erosion.

Third, our dataset choices are intentionally narrow. While GSM8K is a clean and benign dataset that is well-suited for isolating catastrophic forgetting, it represents only a single domain. Other benign domains may interact differently with safety alignment. An analagous statement can be made for the BeaverTails dataset.

Fourth, we evaluate LoX at a single extrapolation coefficient, as specified in the original paper. It's entirely possible that different coefficients yield meaninfully different robustness results under our sequential attack.

Finally, the combined use of ASR and utility metrics is better than either on their own, but remains limited. ASR, even with a strong judge model like Llama-3-8B-Guard, may miss subtle policy violations or fail to capture the severity of unsafe outputs. Utility scores similarly provide only a coarse measure of capability preservation. Furthermore, the choice of judge model may itself meaningfully affect ASR measurements. 

% Our project is a focused test of whether LoX still works under a two-stage sequential fine-tuning attack, so there are a few clear limitations. First, we only evaluate one defense (LoX) and one main fine-tuning setup. We note that catastrophic forgetting and how much the model forgets towards unsafe outputs varies based on different training choices like learning rate, number of steps, LoRA rank, or even a different PEFT method.

% Second, our results are being limited to the specific model we choose to test. Safety robustness can vary greatly between model sizes as well as families, so findings on one specific model may not transfer directly to others. Third, our dataset choices are narrow. While GSM8K is a good and clean benign dataset for isolating if forgetting occurs, it also is only math reasoning which means that other benign areas might interact differently with safety alignment. Furthermore, on the harmful side, BeaverTails covers most unsafe categories, but it is still just one dataset and may not represent every real-world up to date jailbreak style or harmful prompt injection.

% Lastly, our evaluation uses attack success rate (ASR) as our primary safety metric, which is useful for what we want to achieve but not perfect. ASR scoring can miss subtle policy violations or fail to capture how severe an unsafe response is. An increase in refusals could reflect over-refusal, which is a different kind of failure that our metrics do not properly capture.

\paragraph{Future Work.}
An immediate extension of this work is to broaden the evaluation across defenses, models, and datsets. On the defense side, comparing LoX against a more complete spectrum of mitigation strategies, such as alignment-stage defenses, fine-tuning-stage constraints like adversarial training, and post-hoc safety recovery methods, would clarify whether sequential catastrophic forgetting represents a general blind spot across defense types or a vulnerability specific to post-alignment approaches like LoX.

On the model side, testing across multiple open-weight architectures and parameter scales would help determine whether sequential forgetting is a universal property of safety-aligned models or one that is specific to certain architectures or training paradigms. Larger models, or those aligned using DPO rather than RLHF, may express meaningfully different vulnerabilities.

Broadening the dataset choices (on both stages of the attack) would also help generalize our conclusions. On the benign side, domains such as code generation or instruction following might reveal whether different semantics produce different rates of alignment erosion. On the harmful side, incorporating additional benchmarks beyond BeaverTails wouild reduce the risk that our conclusions are specific to that dataset's particular distribution of unsafe content.

Finally, and most ideally, future work would develop practical solutions to the vulnerabilities presented by this paper.

\section*{Statement of Work}

This section should clearly describe the contributions of each group member. This section does not count towards the 6 page limit.

\begin{itemize}
    \item \textbf{Jorge Manuel Torre:} Proposed work, Dataset and Models. Started setting up shared directory for project environment.
    \item \textbf{Elias Mengistu:} Evaluation Metrics, Benign Fine Tuning Dataset.
    \item \textbf{Alex Levy:} Abstract. Problem Statement \& Threat Model, Related Work, Proposed Work, Limitations and Future Work.
    \item \textbf{Jeevan Shahi:} The presentation slides on limitations and future work, abstract, and idea generation for project. Milestone 2 slides as well as some of the documentation associated.
\end{itemize}

\section*{Usage of AI}

AI was used here to make the grammar and layout of this proposal more professional. Specifically, we gave AI our rough draft and implemented its improvements and suggestions to produce the final paper. We used the Chat-GPT5 and Claude Sonnet \& Opus 4.6 models.

{\footnotesize \bibliographystyle{acm}
\bibliography{sample}}


\theendnotes

Our work so far has centered around recreating the original results of the LoX paper based on the artifacts linked in their results. We've run into limitations with memory, gaining access to the models on HuggingFace, and getting the system running on the glogin cluster.

\end{document}